\section{Implementation and Deployment}
\label{sec:implementation}

\subsection{Code organisation}

Each sub-problem of \Cref{sec:architecture} maps to one place in the tree, which
is the practical form of the rule that a component owns one decision.

\begin{table}[H]
  \centering\small
  \setlength{\tabcolsep}{5pt}
  \begin{tabularx}{\textwidth}{@{}>{\hsize=0.46\hsize}Y>{\hsize=0.92\hsize}Y>{\hsize=1.62\hsize}Y@{}}
    \toprule
    Sub-problem & Module & Responsibility \\
    \midrule
    S1 Knowledge acquisition & \modpath{ingestion/} & Document loading, multi-sample extraction with a judge model, and the canonical interface representation shared by both platforms \\
    \addlinespace
    S2 Test selection & \modpath{planner/} & Both planner modes, the graph tools, the validating proposal gate, coverage and the exploration directive, and the pluggable source registry \\
    \addlinespace
    S3 Execution & \modpath{clients/executor_runner.py}, \modpath{web_player/} & Android and web execution, each with its own observation, action and failure-classification code, sharing one attribution vocabulary \\
    \addlinespace
    S4 Evidence reduction & \modpath{rag_api/findings.py} & Trajectory to typed, deduplicated findings with an independent-observation count and a lifecycle \\
    \addlinespace
    S5 Persistence and retrieval & \modpath{rag_api/} & The sole database owner: schema, constraints, vector indexes, retrieval, the state-identity cascade, learning, risk, metrics and anomalies \\
    \addlinespace
    S6 Operation & \modpath{gateway/}, \modpath{dashboard-react/}, \modpath{targets/} & The single entry point, the operator dashboard, the exportable run report, and the validated target profiles \\
    \bottomrule
  \end{tabularx}
\end{table}

\subsection{How the design is verified}

Three mechanisms verify the design without consuming a device or a model
budget, which matters because the alternative is to discover a regression only
during a campaign that costs seven hours of device time.

A \emph{simulator runner} exercises the full planning and logging loop with a
stubbed player, so graph writes, the coverage calculation and the attribution
path can be tested without a handset. A set of \emph{configuration guards}
asserts the invariants the design depends on, including that the attribution
sets are disjoint and that no application-specific vocabulary has entered the
code. And the \emph{state-identity} and \emph{observation} modules are pure
functions over recorded inputs, so the cascade is tested against fixed key sets
rather than against a live device.

The repository carries an automated suite of six modules behind a single entry
point. At the time of writing 124 checks pass with no external service, the
findings module additionally requiring a live database. The suite is weighted
deliberately toward the parts that are pure functions over recorded inputs,
because those are the parts where a regression would be silent.

\subsection{Runtime topology}

\begin{figure}[H]
  \centering
  \includegraphics[width=0.92\textwidth]{12-runtime-topology.png}
  \caption[Runtime topology]{Four processes on one host, one device over the
  Android Debug Bridge, and exactly one external dependency.}
  \label{fig:topology}
\end{figure}

A launch script starts the retrieval service and the gateway under Uvicorn,
waits for both health endpoints, and optionally starts the campaign loop. Neo4j
runs locally. The Android device is a physical handset or an emulator attached
over the Android Debug Bridge, carrying the companion application that provides
the accessibility service and the input method. Model inference is the only
component reached over the network.

\subsection{Configuration, and adding a target}

Configuration is layered deliberately. Process-wide defaults live in a settings
module read once at import, which is why a campaign is launched as a subprocess
rather than in-process: the settings object is global, so running a second
target inside a live server would leak that target's configuration into every
other request.

Per-application configuration is a validated JSON target profile. The dashboard
can list, validate, save, ingest and run one, so adding an application is a
data change rather than a code change. This is the mechanism behind the claim
that five targets were driven without modifying the system, and it is what a
QA engineer would use to point the agent at something new.

Secrets live only in an environment file excluded from version control. The
test-account password reaches the player and never the planner.

\subsection{Operating a campaign}

A campaign is operated from the dashboard: the live session, the growing
application model, both log streams and the planner trace are visible while it
runs, and a run report exports as a PDF afterwards. Three artefact streams are
written to disk for inspection after the fact: structured application logs, one
trajectory directory per execution, and the degradation register.

The graph is the durable record and is the thing that must be backed up.
Snapshots are held outside version control under the same controls as the test
credentials.

\subsection{Scope of the deployment claim}

This is a local, developer-operated deployment on a trusted host, and that is
stated as a boundary rather than discovered later as a defect. The services
bind to a host-local interface by default and support key-based authentication.
Deploying beyond a single trusted host would require uniform authentication
across every diagnostic handler and containment checks on the paths that accept
a target profile. Neither has been done, because neither was needed for the
deployment this system was built for.
